\section{2026 年 11 月核验增补}

本节按论文首次提交时间排列。每项方法附有可核验的外部入选依据。

% METHOD: 2607.00746
\paragraph{GaussianFusion：让相机与 LiDAR 在连续高斯表示中交换特征}
2026.07~\cite{2607_00746}。GaussianFusion 在共享的三维高斯中编码相机和 LiDAR 特征，并迭代更新原语的属性。统一场景表示可送入相应的 BEV 检测头或占据预测头，使两项任务共享融合结果。

在 nuScenes 验证集上，论文的多模态目标检测报告 74.0 NDS 与 71.7 mAP，单帧语义占据设置报告 28.65 mIoU。两类分数对应不同的任务指标。作者指出把多帧信息简单对齐仍未充分建模运动。作者单位包括腾讯自动驾驶实验室。

\textbf{入选依据} 论文署名单位包括 Tencent（\href{https://arxiv.org/html/2607.00746v1}{论文原文}）。

% METHOD: 2607.14203
\paragraph{Instant NuRec：从短驾驶日志前馈生成分层高斯场景}
2026.07~\cite{2607_14203}。Instant NuRec 输入短的多相机驾驶图像序列，通过共享编码器与若干解码头一次预测可渲染的分层高斯，用以将日志重建接到驾驶仿真。分层表达让不同部分分别控制，而不必逐场景从头优化完整场景。

论文把公共 Waymo 的新视图评价与内部大规模数据的重建和闭环仿真分开，并以消融检查关键设计。作者指出更少的原语可能漏掉细结构，显著不同于训练条件的相机布局限制迁移。作者署名为 NVIDIA，论文链接该机构的公开实现。

\textbf{入选依据} 论文署名单位包括 NVIDIA（\href{https://arxiv.org/html/2607.14203v2}{论文原文}）。

% METHOD: 2609.06948
\paragraph{PRG-Fusion：按重建证据决定保留、修复或生成}
2026.09~\cite{2609_06948}。沿日志以外的轨迹渲染驾驶视频时，同一图像可能同时包含可信的重建表面、受退化影响但仍有测量支撑的区域，以及完全缺少观测的区域。PRG-Fusion 根据三维高斯渲染和几何证据按像素划分这三类区域。可信部分保留高斯外观，中间部分用 LiDAR 投影约束结构，缺少支撑的部分交给视频先验补全。

作者在 Waymo 验证集选取八个含较多动态交通参与者的片段进行评价。方法依赖预先重建与对齐的 LiDAR；比较轨迹外推时，应分别报告可由重建或测量支持的区域，以及由生成模型补齐的区域。

\textbf{入选依据} 论文署名单位包括 Yootta（\href{https://arxiv.org/abs/2609.06948}{论文原文}）。

% METHOD: 2609.20756
\paragraph{OPTED：把昂贵的渲染留在端到端驾驶学生的监督阶段}
2026.09~\cite{2609_20756}。OPTED 让能获取特权状态的教师在不渲染图像的向量模拟器中从奖励学习，在基于重建实录的闭环模拟器中为视觉学生提供行为监督。训练流程把强化学习探索与高成本传感器渲染分开，高斯重建在相关实验里充当视觉模拟环境。

作者用 AlpaSim 和 WOMD 不同协议比较驾驶表现。AlpaSim 中的非反应式日志轨迹限制了碰撞解释，教师与视觉学生的可观测性也不同。论文列有 NVIDIA Research 署名，该条属于高斯仿真辅助驾驶学习。

\textbf{入选依据} 论文署名单位包括 NVIDIA（\href{https://arxiv.org/html/2609.20756v1}{论文原文}）。作为 3DGS 支持闭环仿真的应用研究，主要贡献是策略蒸馏。

% METHOD: 2609.39841
\paragraph{DyRAD：分离动态雷达反射点与传感器点扩散}
2026.09~\cite{2609_39841}。DyRAD 将雷达距离、方位、多普勒张量的响应分为静态或由轨迹驱动的点反射体，以及依赖雷达配置的固定点扩散函数。传感器响应的宽展不必被写进场景结构，已有结构因此可在论文考察的不同处理配置下再次渲染。

论文比较真实 RADIal 和 Boreas、合成轨迹偏移，并区分真实数据的循环一致性和合成偏移视角的直接真值。现有实现基于地平面，反射体对应物体的初始指派固定，尚未完整建模多径。官方仓库有训练与评测入口。

\textbf{入选依据} \href{https://github.com/dyrad-nvs/dyrad}{公开实现}已找到与论文对应的代码文件。
